\documentclass[10pt,a4paper]{amsart}
\usepackage[margin=19mm]{geometry}
\usepackage{amsmath,amssymb,mathtools}
\usepackage[scr=rsfs,cal=euler]{mathalfa}
\usepackage{xfrac}
\usepackage[unicode,hidelinks,bookmarks=false]{hyperref}
\newcommand{\ie}{i.e.\ }
\newcommand{\cf}{cf.\ }
\newcommand{\resp}{respectively\ }
\newcommand{\Subsection}{\subsection}
\providecommand{\nobreakdash}{\nobreak\hbox{-}}
\setlength{\emergencystretch}{2em}
\multlinegap0pt
\usepackage{amssymb}
\usepackage{mathtools}
\newtheorem{lemma}{Lemma}
\newcommand{\inj}{\operatorname{inj}}
\title{\fontsize{20}{24}\selectfont \bfseries Tur\'an-good monotonicity thresholds}
\author{Yuanpei Wang\textsuperscript{*} \qquad Liying Kang\textsuperscript{\(\dagger\)} \qquad Xiamiao Zhao\textsuperscript{\(\ddagger\)}}
\date{Source: arXiv:2608.17134v2 (CC BY 4.0)}
\begin{document}
\maketitle

\noindent\emph{Source and licence.} \fontsize{20}{24}\selectfont \bfseries Tur\'an-good monotonicity thresholds. Version arXiv:2608.17134v2 (CC BY 4.0), distributed by arXiv under the Creative Commons Attribution 4.0 International licence, http://creativecommons.org/licenses/by/4.0/, as recorded in the arXiv licence metadata for this exact version and shown on the abstract page https://arxiv.org/abs/2608.17134v2. Full licence notice: \texttt{licenses/arxiv-2608.17134-CC-BY-4.0.txt}.\par\smallskip

\noindent\textit{Editorial scope.} Counted unit: the article's lemma lem:balancing (source0.tex lines 647--661). The article's complete original proof of it is delivered with it (source0.tex lines 663--769, one \texttt{proof} environment of 107 lines). No mathematics is written, repaired or re-derived: the statement and the proof are the article's own lines, in the article's own notation, and no retained line is substituted or re-flowed.  Retained uncounted as the source-local context the proof needs: lines 548--555.  Nothing was omitted from any retained span, so the package carries no omission record.  No retained line contains a citation, so the wrapper carries no bibliography and no bibliography entry.  No retained line contains a figure environment, an image-inclusion command or drawing code, so no image is delivered and the package declares no asset dependency.  Editorial formatting only, in addition: an AMS \texttt{amsart} preamble that restates the article's own declarations and macros used by the retained lines, namely the environments \texttt{lemma} on separate counters, together with the standard packages those lines need.  The retained environments print in this document as Lemma 1 in order of appearance, whereas the article prints the same items with the numbering of its own full document.  The complete native source of the article as distributed in the arXiv e-print for this exact version is bundled unchanged as \texttt{sources/207873/source0.tex}.
\begin{lemma}
\label{lem:two-part-finite-difference}
Let $F$ be a graph with $t:=e(F)\ge 1$. Let $a,b$ be nonnegative integers with $a\ge b+2$, and define $\ell:=a+b$ and $d:=a-b-1$. Then
\[
\left|\inj(F,K_{a-1,b+1})-\inj(F,K_{a,b})\right|
\le 4t^2d\ell^{v(F)-2}.
\]
\end{lemma}

\begin{lemma}
\label{lem:balancing}
Let $K=K_{a_1,\ldots,a_r}$ have parts $A_1,\ldots,A_r$, where $a_k:=|A_k|$. Choose $i\ne j$ with $a_i\ge a_j+2$, and let $K'$ be obtained by moving one vertex from $A_i$ to $A_j$. Define
\[
d:=a_i-a_j-1=e(K')-e(K),\qquad
\ell:=a_i+a_j,\qquad
s:=\frac{\ell}{n},\qquad
\sigma_2:=\sum_{k=1}^r\left(\frac{a_k}{n}\right)^2.
\]
Then
\[
\inj(H,K')-\inj(H,K)
\ge 2mdn^{h-2}\left(1-5\frac{J}{m}s-5ms^2-m\sigma_2-O_H(1/n)\right).
\]
\end{lemma}

\begin{proof}
For $W\subseteq V(H)$, let $R(W)$ be the number of injective homomorphisms from $H-W$ to the complete $(r-2)$-partite graph induced by the parts other than $A_i,A_j$. Splitting each embedding according to $W=\varphi^{-1}(A_i\cup A_j)$ gives
\begin{align*}
\inj(H,K)
&=\sum_{W\subseteq V(H)}\inj(H[W],K_{a_i,a_j})R(W),\\
\inj(H,K')
&=\sum_{W\subseteq V(H)}\inj(H[W],K_{a_i-1,a_j+1})R(W).
\end{align*}
For each $W\subseteq V(H)$, define
\[
\Delta_W:=\inj(H[W],K_{a_i-1,a_j+1})-\inj(H[W],K_{a_i,a_j}).
\]
It follows that
\[
\inj(H,K')-\inj(H,K)=\sum_{W\subseteq V(H)}\Delta_WR(W).
\]
If $e(H[W])=0$, both embedding counts in the definition of $\Delta_W$ equal $(\ell)_{|W|}$, so $\Delta_W=0$. Define
\[
\Delta^{(1)}
:=\sum_{\substack{W\subseteq V(H)\\e(H[W])=1}}\Delta_WR(W),
\qquad
\Delta^{(\ge2)}
:=\sum_{\substack{W\subseteq V(H)\\e(H[W])\ge2}}\Delta_WR(W).
\]
Then
\[
\inj(H,K')-\inj(H,K)=\Delta^{(1)}+\Delta^{(\ge2)}.
\]

We first estimate $\Delta^{(1)}$. If $e(H[W])=1$, then $H[W]$ consists of one edge and isolated vertices. In this case
\[
\inj(H[W],K_{a_i,a_j})
=2a_ia_j(\ell-2)_{|W|-2}.
\]
Since $2\bigl((a_i-1)(a_j+1)-a_ia_j\bigr)=2d$, we have
\begin{equation}
\label{eq:one-edge-delta}
\Delta_W=2d(\ell-2)_{|W|-2}.
\end{equation}

Fix $uv\in E(H)$. To sum \eqref{eq:one-edge-delta} over the sets $W$ satisfying $E(H[W])=\{uv\}$, regard the images of $u$ and $v$ as fixed in the two selected parts and map the remaining $h-2$ vertices by a uniformly random injection into the remaining $n-2$ host vertices. Let $S$ be the set of the remaining $\ell-2$ vertices in $A_i\cup A_j$. Such an injection contributes to the desired sum provided none of the following events occurs.
\begin{enumerate}
\renewcommand{\labelenumi}{(\roman{enumi})}
\item Some vertex in $(N_H(u)\cup N_H(v))\setminus\{u,v\}$ is mapped into $S$. There are at most $d_H(u)+d_H(v)-2$ such vertices, and each is mapped into $S$ with probability $(\ell-2)/(n-2)=s+O(1/n)$. Hence this event has probability at most $\bigl(d_H(u)+d_H(v)-2\bigr)(s+O(1/n))$.
\item For some edge $ab\in E(H)$ vertex-disjoint from $uv$, both $a$ and $b$ are mapped into $S$. Then $a,b\in W$, so $E(H[W])\ne\{uv\}$. For each of at most $m$ such edges, the probability is $(\ell-2)_2/(n-2)_2=s^2+O(1/n)$, so the total probability is at most $m(s^2+O(1/n))$.
\item For some $k\notin\{i,j\}$, both endpoints of an edge of $H$ are mapped into $A_k$. For each edge, the probability is at most $\sum_{k\notin\{i,j\}}(a_k)_2/(n-2)_2\le \sigma_2+O(1/n)$. Summing over the $m$ edges gives $m(\sigma_2+O(1/n))$.
\end{enumerate}
By the union bound, the probability that at least one of these events occurs is at most
\[
\bigl(d_H(u)+d_H(v)-2\bigr)s+ms^2+m\sigma_2+O_H(1/n).
\]
Hence the terms corresponding to the sets $W$ satisfying $E(H[W])=\{uv\}$ sum to at least
\[
2d(n-2)_{h-2}
\left(1-\bigl(d_H(u)+d_H(v)-2\bigr)s-ms^2-m\sigma_2-O_H(1/n)\right).
\]
Summing over $uv\in E(H)$ and using
\[
\sum_{uv\in E(H)}\bigl(d_H(u)+d_H(v)-2\bigr)=J,
\qquad
(n-2)_{h-2}=n^{h-2}(1+O_H(1/n)),
\]
we obtain
\begin{equation}
\label{eq:balancing-one-edge}
\Delta^{(1)}
\ge 2mdn^{h-2}\left(1-\frac{J}{m}s-ms^2-m\sigma_2-O_H(1/n)\right).
\end{equation}

For $e(H[W])\ge2$, Lemma \ref{lem:two-part-finite-difference} gives $|\Delta_W|\le 4e(H[W])^2d\ell^{|W|-2}$. Therefore,
\[
\Delta^{(\ge2)}
\ge -4d\sum_{\substack{W\subseteq V(H)\\e(H[W])\ge2}}
e(H[W])^2\ell^{|W|-2}R(W).
\]
Since $R(W)\le(n-\ell)_{h-|W|}\le(n-\ell)^{h-|W|}$ and $s=\ell/n$,
it follows that
\[
\Delta^{(\ge2)}
\ge -4dn^{h-2}
\sum_{\substack{W\subseteq V(H)\\e(H[W])\ge2}}
e(H[W])^2s^{|W|-2}(1-s)^{h-|W|}.
\]
When $e(H[W])\ge2$, $e(H[W])^2\le2e(H[W])(e(H[W])-1)$, where $e(H[W])(e(H[W])-1)$ counts the ordered pairs of distinct edges in $H[W]$. For an ordered pair $(e,f)$ of distinct edges of $H$, let $U_{e,f}:=V(e)\cup V(f)$. Hence
\[
4\sum_{\substack{W\subseteq V(H)\\e(H[W])\ge2}}e(H[W])^2s^{|W|-2}(1-s)^{h-|W|}
\le 8\sum_{\substack{(e,f)\in E(H)^2\\e\ne f}}\ \sum_{U_{e,f}\subseteq W\subseteq V(H)}s^{|W|-2}(1-s)^{h-|W|}.
\]
For each fixed pair $(e,f)$, every set $W$ satisfying $U_{e,f}\subseteq W\subseteq V(H)$ is uniquely of the form $W=U_{e,f}\cup Z$, where $Z\subseteq V(H)\setminus U_{e,f}$. Hence
\[
\sum_{U_{e,f}\subseteq W\subseteq V(H)}s^{|W|-2}(1-s)^{h-|W|}
=s^{|U_{e,f}|-2}\sum_{Z\subseteq V(H)\setminus U_{e,f}}s^{|Z|}(1-s)^{h-|U_{e,f}|-|Z|}
=s^{|U_{e,f}|-2}.
\]
If $e$ and $f$ share one endpoint, then $|U_{e,f}|=3$; otherwise $|U_{e,f}|=4$. There are $J$ ordered pairs of the first type and at most $m^2$ of the second. Therefore,
\begin{equation}
\label{eq:balancing-multiple-edge}
\Delta^{(\ge2)}\ge -8dn^{h-2}(Js+m^2s^2).
\end{equation}
Combining \eqref{eq:balancing-one-edge} and \eqref{eq:balancing-multiple-edge} gives
\[
\inj(H,K')-\inj(H,K)
\ge2mdn^{h-2}
\left(1-5\frac{J}{m}s-5ms^2-m\sigma_2-O_H(1/n)\right),
\]
as required.
\end{proof}

\end{document}
